% \chapter*{Conclusões e Trabalhos Futuros}\label{cap:conclusao}
\chapter{Conclusões e Próximos Passos}\label{cap:conclusao}
% \addcontentsline{toc}{chapter}{Conclusão e Trabalhos Futuros}

\section{Conclusão}

Este trabalho apresentou o desenvolvimento inicial de uma abordagem baseada em computação quântica de variáveis contínuas para a resolução de problemas de otimização combinatória, com ênfase no Problema do Caixeiro Viajante como caso de estudo. Duas formulações Hamiltonianas foram propostas e implementadas: a primeira baseada em um mapeamento direto entre variáveis binárias e operadores de posição com comportamento binário imposto energeticamente, e a segunda fundamentada em variáveis inteiras representadas por estados de Fock truncados, na qual cada vértice do grafo é associado a um qumode cujo nível de ocupação codifica sua posição na rota.

Os resultados obtidos demonstram a consistência das formulações propostas. Em todas as instâncias avaliadas, com $N \in \{3, 4, 5, 6, 7, 8\}$ cidades, os estados de menor energia identificados pelo algoritmo quântico coincidiram com os caminhos de menor custo encontrados pelo algoritmo clássico de força bruta, validando a correção do formalismo desenvolvido. Adicionalmente, a formulação apresenta uma vantagem estrutural relevante em relação às abordagens baseadas em qubits: um grafo de $N$ vértices pode ser representado com apenas $N$ qumodes com corte $d = N$, enquanto formulações baseadas em qubits, como QUBO ou Ising com codificação \textit{one-hot}\footnote{A codificação \textit{one-hot} representa uma escolha discreta por um vetor binário no qual exatamente uma posição tem valor $1$, indicando a alternativa selecionada, enquanto todas as demais têm valor $0$.}, requerem $N^2$ variáveis binárias, resultando em espaços de estados significativamente maiores para uma mesma instância, conforme ilustrado na \autoref{tab:comparacao}.

Embora os experimentos estejam restritos a instâncias de pequena dimensão em razão do elevado custo computacional associado à simulação clássica de sistemas de variáveis contínuas, os resultados preliminares obtidos estabelecem uma base sólida para as etapas seguintes do trabalho, descritas na seção a seguir.

\section{Próximos Passos}\label{sec:next_steps}

A partir dos resultados obtidos, o desenvolvimento do trabalho 
estrutura-se em duas frentes complementares: a aplicação de algoritmos de otimização quântica variacional sobre os Hamiltonianos já construídos, e a generalização do formalismo proposto para variantes mais complexas do problema de roteamento de veículos.

Os Hamiltonianos construídos serão empregados como funções de custo em modelos de otimização quântica variacional, com destaque para o VQE e o QAOA, apresentados na \autoref{sec:quantum_approach} e aqui adaptados ao contexto dos Hamiltonianos formulados em variáveis contínuas. Nessas abordagens, estados parametrizados serão preparados por circuitos variacionais compostos por portas gaussianas e, quando necessário, operações não gaussianas, sendo então avaliados iterativamente. Os parâmetros variacionais serão ajustados por rotinas de otimização clássica baseadas em gradiente ou métodos livres de derivadas, a depender da estrutura do Hamiltoniano e do comportamento numérico observado.

A avaliação de desempenho não se restringirá aos modelos quânticos. Para fins de comparação, serão empregados métodos clássicos de referência, incluindo algoritmos exatos em instâncias de pequena dimensão. Essa abordagem permitirá analisar, de forma sistemática, a qualidade das soluções encontradas e o custo computacional envolvido, bem como investigar a estabilidade dos métodos frente à escolha de hiperparâmetros. Os resultados obtidos servirão de base para uma análise comparativa entre as abordagens clássicas e quânticas, permitindo avaliar o potencial e as limitações de cada estratégia aplicada aos Hamiltonianos construídos.

Como extensão imediata do formalismo para o caso de múltiplos veículos, propõe-se associar a cada vértice $V_k$ não apenas a variável inteira $n_k \in \{1, 2, \ldots, N\}$, que indica sua posição de visita, mas também uma variável discreta $v_k \in \{1, 2, \ldots, M\}$, que identifica o veículo responsável por seu atendimento, sendo $M$ o número total de veículos disponíveis. Desse modo, cada configuração passa a ser descrita pelo estado:

\begin{equation}
\ket{(n_1, v_1),(n_2, v_2), \ldots,(n_N, v_N)},
\end{equation}

\noindent em que o par $(n_k, v_k)$ informa, simultaneamente, a ordem de visita do vértice $V_k$ e o veículo ao qual ele foi atribuído. Nessa representação, é conveniente introduzir o projetor:

\begin{equation}
    \Pi^{(a,\mu)}_k = \ketbra{a}{a}_k^{(n)}\otimes \ketbra{\mu}{\mu}_k^{(v)},
\end{equation}

\noindent que seleciona os estados nos quais o vértice $V_k$ ocupa a posição $a$ na rota do veículo $\mu$. Com essa nova codificação, o Hamiltoniano de custo pode ser generalizado de modo a contabilizar apenas deslocamentos entre vértices consecutivos atendidos pelo mesmo veículo:

\begin{equation}
H^{(M)}_{cost} = \sum_{\mu=1}^{M} \sum_{a=1}^{N-1} \sum_{i=1}^{N} 
\sum_{j=1}^{N} d_{ij} \left(|a\rangle \langle a|^{(n)}_k \otimes 
|\mu\rangle \langle \mu|^{(v)}_k\right) \left(|a\rangle \langle 
a+1|^{(n)}_j \otimes |\mu\rangle \langle \mu|^{(v)}_j\right),
\end{equation}

\noindent onde $d_{ij}$ representa o custo de deslocamento entre os vértices $V_i$ e $V_j$. Para assegurar a consistência combinatória da solução, define-se o operador de ocupação

\begin{equation}
\hat{N}_{a,\mu} = \sum_{k=1}^{N} \Pi^{(a,\mu)}_k,
\end{equation}

\noindent a partir do qual se introduz o termo de penalização

\begin{equation}
H_{occ} = \sum_{\mu=1}^{M} \sum_{a=1}^{N} \hat{N}_{a,\mu} 
\left(\hat{N}_{a,\mu} - 1\right),
\end{equation}

\noindent que penaliza configurações em que mais de um cliente é associado à mesma posição $a$ do veículo $\mu$. Para desencorajar rotas com posições vazias intercaladas entre posições ocupadas, acrescenta-se

\begin{equation}
H_{gap} = \sum_{\mu=1}^{M} \sum_{a=1}^{N-1} \left(1 - 
\hat{N}_{a,\mu}\right)\hat{N}_{a+1,\mu},
\end{equation}

\noindent de modo que o Hamiltoniano total assuma a forma

\begin{equation}
H_{VRP} = H^{(M)}_{cost} + \lambda_1 H_{occ} + \lambda_2 H_{gap},
\end{equation}

\noindent com $\lambda_1, \lambda_2 > 0$ parâmetros de penalização.

Como etapa posterior, pode-se incorporar ao formalismo uma restrição de capacidade. Se $q_k$ denota a demanda do cliente $V_k$ e $Q_\mu$ a capacidade do veículo $\mu$, define-se o operador

\begin{equation}
\hat{y}_{k,\mu} = \sum_{a=1}^{N} \Pi^{(a,\mu)}_k,
\end{equation}

\noindent que projeta sobre os estados em que o cliente $V_k$ é atendido pelo veículo $\mu$, independentemente de sua posição na rota. A carga total associada ao veículo $\mu$ pode então ser escrita como

\begin{equation}
\hat{Q}_\mu = \sum_{k=1}^{N} q_k \hat{y}_{k,\mu}.
\end{equation}

Com isso, um possível Hamiltoniano de penalização para a restrição de capacidade é dado por

\begin{equation}
H_{cap} = \sum_{\mu=1}^{M} \Theta\left(\hat{Q}_\mu - Q_\mu\right) 
\left(\hat{Q}_\mu - Q_\mu\right)^2,
\end{equation}

\noindent em que $\Theta(\cdot)$ representa a função de Heaviside, atribuindo energia positiva apenas às configurações em que a demanda total dos clientes atribuídos ao veículo $\mu$ excede sua capacidade máxima.

De modo análogo, a inclusão de janelas de tempo pode ser proposta mediante a introdução de operadores de tempo de chegada ao longo de cada rota. Se $[A_i, B_i]$ representa a janela de atendimento do cliente $V_i$, define-se $\hat{T}_{a,\mu}$ como o operador associado ao instante de chegada do veículo $\mu$ ao cliente visitado na posição $a$. A consistência temporal da rota pode ser imposta por meio do termo

\begin{equation}
H_{temp} = \sum_{\mu=1}^{M} \sum_{a=1}^{N-1} \left(\hat{T}_{a+1,\mu} 
- \hat{T}_{a,\mu} - \sum_{i=1}^{N}\sum_{j=1}^{N}(\tau_{ij} + s_i) 
\Pi^{(a,\mu)}_i \Pi^{(a+1,\mu)}_j\right)^2,
\end{equation}

\noindent em que $\tau_{ij}$ denota o tempo de deslocamento entre os clientes $V_i$ e $V_j$, e $s_i$ o tempo de serviço em $V_i$. A penalização associada às janelas de tempo pode então ser escrita como

\begin{equation}
H_{jt} = \sum_{\mu=1}^{M} \sum_{a=1}^{N} \sum_{i=1}^{N} 
\Pi^{(a,\mu)}_i \left[\Theta\left(A_i - \hat{T}_{a,\mu}\right) 
\left(A_i - \hat{T}_{a,\mu}\right)^2 + \Theta\left(\hat{T}_{a,\mu} 
- B_i\right)\left(\hat{T}_{a,\mu} - B_i\right)^2\right],
\end{equation}

\noindent em que a energia do sistema cresce sempre que o atendimento ocorre antes da abertura ou após o fechamento da janela admissível.

O cronograma de desenvolvimento previsto para o período posterior à qualificação, organizado ao longo do segundo semestre de 2026, está sintetizado na Tabela~\ref{tab:cronograma}.

    \begin{table}[!htb]
    \centering
    \caption{Cronograma de execução das atividades}
    \label{tab:cronograma}
    \resizebox{\textwidth}{!}{
    \begin{tabular}{|p{6.2cm}|c|c|c|c|c|c|c|c|c|c|c|c|c|}
    \hline

    \textbf{Etapa} & \textbf{Jun} & \textbf{Jul} & \textbf{Ago} & \textbf{Set} & \textbf{Out} & \textbf{Nov} & \textbf{Dez} \\  \hline \hline
    Revisão bibliográfica e ampliação do referencial teórico &  \cellcolor{gray}  &  \cellcolor{gray}  &  \cellcolor{gray}  &     &     &     & \\  \hline
    Desenvolvimento e extensão do Hamiltoniano               &  \cellcolor{lightgray}  &  \cellcolor{lightgray}  &  \cellcolor{lightgray}  &  \cellcolor{lightgray}  &     &     & \\  \hline
    Implementação computacional  & &  \cellcolor{gray}  &  \cellcolor{gray}  &  \cellcolor{gray}  &  \cellcolor{gray}  &         & \\  \hline
    Simulações, análise de resultados e validação            &     &     &  \cellcolor{lightgray}  &  \cellcolor{lightgray}  &  \cellcolor{lightgray}  &     & \\  \hline
    Redação e consolidação da tese &  \cellcolor{gray}  &  \cellcolor{gray}  &  \cellcolor{gray}  &  \cellcolor{gray}  &  \cellcolor{gray}  &  \cellcolor{gray}  & \\  \hline
    Defesa da Tese &    &    &    &     &     &      & \cellcolor{lightgray}  \\  \hline
    \end{tabular}
    }
    \fonte{Autoria Própria (2026)}
    \end{table}

    